\chapter{Zdrojový kód aplikace}
\label{app:kod}

Kompletní zdrojový kód aplikace Foody, podklady z~obou fází uživatelského testování, předkompilovaný release build pro Android i~doprovodné materiály jsou dostupné v~Git repozitáři Fakulty elektrotechnické ČVUT na adrese:

\begin{center}
\url{https://gitlab.fel.cvut.cz/andrajak/diplomka-foody}
\end{center}

Repozitář obsahuje zdrojový kód aplikace ve frameworku Flutter (\texttt{lib/}), nativní platformové projekty pro Android a~iOS, validační benchmark přesnosti AI odhadu na datasetu Nutrition5k (\texttt{test/ai\_benchmark/}), anonymizované podklady z~uživatelského testování (\texttt{testovani/kratkodobe/} pro účastníky P1--P3 a~\texttt{testovani/dlouhodobe/} pro LP1--LP6) a~předkompilovaný release build aplikace pro Android v~souboru \texttt{foody-release.apk}.

Aplikace byla primárně vyvíjena pro iOS (referenční zařízení iPhone 16 Pro s~operačním systémem iOS 26), na vybraných Android zařízeních byla doplňkově ověřena pro účely uživatelského testování. Návod na sestavení, konfiguraci klíče \texttt{OPENAI\_API\_KEY} v~souboru \texttt{.env} i~odkaz na HiFi prototyp v~prostředí Figma jsou v~souboru \texttt{README.md} v~kořenovém adresáři repozitáře.

\paragraph{Orientace v~repozitáři.}
Pro účely reprodukovatelnosti experimentů popsaných v~kapitole~\ref{ch:testovani} jsou klíčové části kódu lokalizovány následovně:
\begin{itemize}
	\item \textbf{Systémové prompty pro model umělé inteligence} (sekce~\ref{subsec:prompty}) v~souboru \texttt{lib/utils/prompt.dart}, obsahujícím šablony pro všechny scénáře (rozpoznání jídla, analýza cvičení, výpočet nutričních cílů, dotazy v~přirozeném jazyce).
	\item \textbf{Sanitizace vstupu a ochrana proti prompt injection} (sekce~\ref{subsec:promptInjection}) v~souboru \texttt{lib/utils/prompt\_sanitizer.dart}.
	\item \textbf{Pipeline pro umělou inteligenci a confidence gate} (sekce~\ref{subsec:confidence}) v~adresáři \texttt{lib/services/ai\_feature/}.
\end{itemize}

\chapter{HiFi prototyp uživatelského rozhraní}
\label{app:hifi}

HiFi prototyp uživatelského rozhraní aplikace Foody vytvořený v~nástroji Figma je k~dispozici jako statický PDF export v~souboru \texttt{DiplomkaHiFi.pdf} na adrese:

\begin{center}
\url{https://gitlab.fel.cvut.cz/andrajak/diplomka-foody/-/blob/main/DiplomkaHiFi.pdf}
\end{center}

Odkaz na interaktivní zdrojový soubor v~prostředí Figma je uveden v~souboru \texttt{README.md} v~kořenovém adresáři repozitáře.

\chapter{Případy užití}
\label{app:pripadyUziti}

Tato příloha obsahuje plné hlavní scénáře osmi případů užití (UC01--UC08), jejichž souhrnný přehled je uveden v~sekci~\ref{sec:pripadyUziti} hlavního textu.
Pro každý případ užití je popsána charakteristika, primární aktér a~hlavní scénář průběhu interakce.

\paragraph{UC01: Záznam jídla z~fotografie.}
\textbf{Charakteristika:} Uživatel chce rychle zapsat jídlo pomocí fotografie a následně výsledek zkontrolovat.
\textbf{Primární aktér:} uživatel.
Hlavní scénář:

\begin{enumerate}
	\item \emph{Uživatel} otevře funkci pro záznam jídla a pořídí fotografii.
	\item \emph{Systém} spustí rozpoznání, zobrazí návrh položek a odhad množství, případně indikaci nejistoty.
	\item \emph{Uživatel} zkontroluje návrh, upraví položky a množství, případně doplní chybějící položky.
	\item \emph{Uživatel} potvrdí uložení záznamu.
	\item \emph{Systém} uloží záznam a aktualizuje denní přehled.
\end{enumerate}

\paragraph{UC02: Oprava výsledku rozpoznání a opakování rozpoznání.}
\textbf{Charakteristika:} Uživatel chce opravit rozpoznané položky a v~případě potřeby vyvolat rozpoznání znovu bez zakládání nového záznamu.
\textbf{Primární aktér:} uživatel.
Hlavní scénář:

\begin{enumerate}
	\item \emph{Uživatel} otevře detail rozpoznaného jídla v~editační obrazovce.
	\item \emph{Systém} zobrazí seznam položek, množství, jednotky a souhrn živin.
	\item \emph{Uživatel} upraví názvy, množství, odstraní nebo přidá položky.
	\item \emph{Uživatel} zvolí možnost znovu rozpoznat, například po úpravě vstupu nebo doplnění informace.
	\item \emph{Systém} provede nové rozpoznání a aktualizuje návrh položek, zachová možnost ruční kontroly.
	\item \emph{Uživatel} potvrdí finální stav a uloží záznam.
\end{enumerate}

\paragraph{UC03: Ruční přidání záznamu jídla.}
\textbf{Charakteristika:} Uživatel chce vytvořit záznam bez použití fotografie nebo umělé inteligence.
\textbf{Primární aktér:} uživatel.
Hlavní scénář:

\begin{enumerate}
	\item \emph{Uživatel} otevře ruční zápis jídla.
	\item \emph{Uživatel} vyhledá položku v~databázi nebo zvolí položku z~historie.
	\item \emph{Uživatel} zadá množství a jednotku, gramy nebo kusy, případně předvolbu porce.
	\item \emph{Uživatel} potvrdí uložení záznamu.
	\item \emph{Systém} uloží záznam a aktualizuje denní přehled.
\end{enumerate}

\paragraph{UC04: Přidání balené potraviny pomocí čárového kódu.}
\textbf{Charakteristika:} Uživatel chce rychle přidat balenou potravinu načtením kódu.
\textbf{Primární aktér:} uživatel.
Hlavní scénář:

\begin{enumerate}
	\item \emph{Uživatel} otevře čtečku kódů a naskenuje čárový kód produktu.
	\item \emph{Systém} vyhledá produkt v~databázi a zobrazí nalezenou položku.
	\item \emph{Uživatel} zvolí množství a jednotku nebo předvolbu porce.
	\item \emph{Uživatel} potvrdí vložení do denního záznamu.
	\item \emph{Systém} uloží záznam a aktualizuje denní přehled.
\end{enumerate}

\paragraph{UC05: Přidání opakovaného jídla z~oblíbených nebo duplikací.}
\textbf{Charakteristika:} Uživatel chce co nejrychleji přidat jídlo, které se často opakuje.
\textbf{Primární aktér:} uživatel.
Hlavní scénář:

\begin{enumerate}
	\item \emph{Uživatel} otevře seznam oblíbených nebo historii záznamů.
	\item \emph{Systém} nabídne oblíbené položky, nedávné záznamy a našeptávání podle historie.
	\item \emph{Uživatel} vybere položku nebo duplikuje celý záznam jídla.
	\item \emph{Uživatel} případně upraví množství a potvrdí uložení.
	\item \emph{Systém} uloží nový záznam a aktualizuje denní přehled.
\end{enumerate}

\paragraph{UC06: Vaření s~režimem plán versus realita.}
\textbf{Charakteristika:} Uživatel chce při vaření zadat plánované ingredience a po dokončení upravit skutečně použité množství.
\textbf{Primární aktér:} uživatel.
Hlavní scénář:

\begin{enumerate}
	\item \emph{Uživatel} vytvoří záznam receptu nebo vaření a zadá plánované ingredience.
	\item \emph{Systém} průběžně počítá plánované souhrnné živiny receptu.
	\item \emph{Uživatel} po uvaření upraví skutečně použité množství, včetně drobných ingrediencí.
	\item \emph{Systém} přepočítá souhrny a nabídne uložení finální verze.
	\item \emph{Uživatel} uloží jídlo do denního záznamu a případně do vlastních receptů pro opakování.
\end{enumerate}

\paragraph{UC07: Nastavení profilu, cílů a dietních omezení.}
\textbf{Charakteristika:} Uživatel chce nastavit cíle, profil a kontext, který ovlivní interpretaci a doporučení.
\textbf{Primární aktér:} uživatel.
Hlavní scénář:

\begin{enumerate}
	\item \emph{Uživatel} otevře nastavení profilu.
	\item \emph{Uživatel} vyplní nebo upraví základní parametry a nastaví cílové hodnoty.
	\item \emph{Uživatel} nastaví dietní omezení nebo intolerance.
	\item \emph{Systém} uloží nastavení a promítne je do přehledů a upozornění.
\end{enumerate}

\paragraph{UC08: Přehledy, dotazy nad daty a export.}
\textbf{Charakteristika:} Uživatel chce vyhodnotit svůj progres a získat odpovědi nad historií, případně data sdílet.
\textbf{Primární aktér:} uživatel.
Hlavní scénář:

\begin{enumerate}
	\item \emph{Uživatel} otevře přehledy a zvolí období, týden nebo měsíc.
	\item \emph{Systém} zobrazí souhrny a trendy, případně upozorní na překročení limitů.
	\item \emph{Uživatel} položí dotaz nad daty, například porovnání období nebo dotaz na makroživiny.
	\item \emph{Systém} zobrazí odpověď a nabídne proklik na relevantní dny.
	\item \emph{Uživatel} volitelně spustí export dat do tabulkového formátu.
\end{enumerate}

\chapter{Anonymizované reporty z~uživatelských rozhovorů}
\label{app:rozhovory}

Anonymizované zápisy ze čtyř strukturovaných rozhovorů s~respondenty R1--R4 popsaných v~sekci~\ref{sec:uzivatelskaStudie} jsou dostupné v~Git repozitáři práce v~adresáři \texttt{testovani/rozhovory/} na adrese:

\begin{center}
\url{https://gitlab.fel.cvut.cz/andrajak/diplomka-foody/-/tree/main/testovani/rozhovory}
\end{center}

Ve zmíněném adresáři jsou k~dispozici individuální zápisy z~rozhovorů se všemi čtyřmi respondenty doplněné o~osnovu rozhovoru použitou moderátorem.

\chapter{Doprovodné materiály k~uživatelskému testování}
\label{app:testovani}

Doprovodné materiály ke krátkodobému moderovanému testování s~participanty P1--P3 popsanému v~sekci~\ref{sec:metodikaTestovani} jsou dostupné v~Git repozitáři práce v~adresáři \texttt{testovani/kratkodobe/} na adrese:

\begin{center}
\url{https://gitlab.fel.cvut.cz/andrajak/diplomka-foody/-/tree/main/testovani/kratkodobe}
\end{center}

Ve zmíněném adresáři jsou k~dispozici individuální záznamy z~testovacích relací jednotlivých participantů v~souborech \texttt{testovani\_\allowbreak P1.md}, \texttt{testovani\_\allowbreak P2.md} a~\texttt{testovani\_\allowbreak P3.md} doplněné o~souhrnný report \texttt{souhrn\_\allowbreak testovani.md}.
Tento report obsahuje kompletní výsledky plnění úloh, vyhodnocení dotazníků SEQ, SUS a~UEQ-S, naměřené hodnoty nefunkčních požadavků, identifikované problémy a~návrhy na zlepšení.

\chapter{Doprovodné materiály k~dlouhodobému testování}
\label{app:dlouhodobeTestovani}

Doprovodné materiály k~dlouhodobému uživatelskému testování v~reálném prostředí s~účastníky LP1--LP6 popsanému v~sekci~\ref{sec:dlouhodobeTestovani} jsou dostupné v~Git repozitáři práce v~adresáři \texttt{testovani/dlouhodobe/} na adrese:

\begin{center}
\url{https://gitlab.fel.cvut.cz/andrajak/diplomka-foody/-/tree/main/testovani/dlouhodobe}
\end{center}

V~podadresáři \texttt{testovani/dlouhodobe/dotazniky/} jsou k~dispozici anonymizované výstupní dotazníky všech šesti účastníků (SUS, Raw NASA-TLX a~otevřené otázky) v~souborech \texttt{vystupni\_\allowbreak dotaznik\_\allowbreak LP1.txt} až \texttt{vystupni\_\allowbreak dotaznik\_\allowbreak LP6.txt} doplněné o~souhrnný report \texttt{souhrn\_\allowbreak dotazniku.md}. V~adresáři \texttt{testovani/dlouhodobe/reporty/} jsou pro každého z~účastníků uloženy úplné exporty behaviorálních dat z~aplikace ve formátu CSV.

\chapter{Použité nástroje umělé inteligence}
\label{app:aiNastroje}

Tato příloha uvádí přehled nástrojů umělé inteligence použitých při přípravě této diplomové práce v~souladu s~Rámcovými pravidly používání umělé inteligence na ČVUT pro studijní a~pedagogické účely v~bakalářském a~navazujícím magisterském studiu. Stvrzení správnosti vygenerovaných výstupů a~odpovědnosti za obsah práce je součástí podepsaného prohlášení.

\begin{itemize}\setlength{\itemsep}{2pt}
	\item \textbf{Claude Code}:
		\begin{itemize}\setlength{\itemsep}{1pt}
			\item revize formulací, kontrola stylistické konzistence a~typografické úpravy textu práce.
			\item asistence při implementační části aplikace.
			\item zpracování dat z~AI benchmarků a~behaviorálních dat z~dlouhodobého testování.
		\end{itemize}
	\item \textbf{Figma Make}: návrh podkladů pro HiFi prototyp (sekce~\ref{sec:prototyp}).
	\item \textbf{OpenAI GPT modely}: klíčová komponenta implementovaného řešení pro rozpoznávání jídel, analýzu cvičení a~generování nutričních doporučení (sekce~\ref{sec:integraceAi}).
\end{itemize}

